\documentclass{article}
\usepackage{icmlko}
\usepackage{graphicx}
\usepackage{booktabs}
\usepackage{amsmath}

\title{겹친다고 말할 검정력조차 없었다 \\
  {\large 그런데 바닥에서 격자를 빼자 네 판이 모두 갈렸다 --- 문제는 신호가 아니라 바닥이었다}}

\begin{document}
\maketitle

\begin{abstract}
직전 노트는 「같은 세계의 바닥이 신호 아래로 내려왔다」로 ③ 을 0.5 칸 밀었고,
프레시 티처가 절제로 그 빈틈의 이름을 무너뜨렸다.
이 노트는 그 절제를 \textbf{사전등록하고 다시 쟀다} --- 그리고 사전등록대로 결론을 그대로 적는다.
\claimx{\textbf{[달력제거 $\wedge$ 첫간격제외]에서 두 구간이 겹친다} ---
진짜 $0.2165\,[0.1980, 0.2368]$ 대 N2 $0.1986\,[0.1777, 0.2194]$.}
{안 겹치면 이름까지 옳게 ③ 이 0 을 벗어난다고 미리 적었다. 겹쳤다.
그러므로 \textbf{「이 자로는 순서를 못 잰다」가 확정}되고 ③ 의 첫째 자는 $0.5 \to 0$ 으로 되돌아간다.}
\claimx{🔴 \textbf{그런데 「겹친다」조차 이 표본이 말할 수 있는 문장이 아니다.}}
{관문 신고가 측정 \emph{전에} 자동으로 찍었다 --- 기대 빈틈 $|-0.0214|$ ÷ MDE $0.0403$ $=$ \textbf{0.532 $<$ 1}.
그러므로 정직한 문장은 「빈틈이 0 이다」가 아니라 \textbf{「빈틈이 0.04025 일 미만이다」}는 상한이다.}
\claimx{🔴 \textbf{이 사이클의 진짜 발견은 판정이 아니라 자다.}}
{바닥이 「격자를 모르는 팔」이라 시간 구조가 전혀 없는 세계에서도 $+0.2827$ 로 유의하게 이기고,
그 이득의 \textbf{92.8\%가 발생률 추정}이다.
바닥에 격자를 넣어 다시 재면 \textbf{네 판이 모두 갈린다} --- [둘 다] 판에서도
진짜 $0.01656\,[0.0101, 0.0234]$ 대 N2 $0.00350\,[-0.0016, +0.0082]$ 로 \textbf{귀무가 0 을 문다}.
다만 그 개선의 크기는 \textbf{23.9 분}이고, 이것은 병기이지 주 판정이 아니다(노트 133).}
\end{abstract}

\section{사전등록이 미리 적은 것}

「어느 쪽이든 끝난다」를 도피구 없이 박았다:
[달력제거 $\wedge$ 첫간격제외]에서 \textbf{안 겹치면} 이름까지 옳게 0 탈출,
\textbf{겹치면} 「이 자로는 순서를 못 잰다」 확정.
헤더 시각은 손으로 안 적었다 --- 자리표시자에 \texttt{st\_mtime} 을 읽어 박고 mtime 을 되돌렸다.

\subsection{자기시험 --- 관이 같다는 것을 먼저 보였다}

직전 두 노트의 정본을 \textbf{11 항목 전부 비트로 재현}했다(덮음 $0.904724581677664$ 까지).
그리고 \textbf{티처의 여덟 수를 여덟 개 다 재현}했다(최대 어긋남 $4.5\times10^{-5}$ --- 티처 쪽 반올림).
\textbf{관이 다르지 않다.}

\section{여덟 수 --- 그리고 판정}

\begin{figure}[h]\centering
\includegraphics[width=\linewidth]{figs/fig1.pdf}
\caption{(a) 네 절제 $\times$ 두 세계. (b) 🔴 관문마다 비를 신고했다 --- 주 판정 관문 자신이 검정력 0 이다.
(c) 🔴 바닥에서 격자를 빼자 네 판이 모두 갈렸다(병기).}
\end{figure}

\begin{table}[h]\centering\small
\begin{tabular}{lrrrr}
\toprule
절제 & 진짜 & N2 & 빈틈 & MDE \\
\midrule
원판 & $+0.7545$ & $+0.2827$ & $+0.3122$ & 0.1682 \\
첫간격제외 & $+0.4479$ & $+0.2062$ & $+0.1436$ & 0.1009 \\
달력제거 & $+0.3238$ & $+0.2190$ & $+0.0404$ & 0.0665 \\
\textbf{둘 다}(주 판정) & $\mathbf{+0.2165}$ & $\mathbf{+0.1986}$ & $\mathbf{-0.0214}$ & \textbf{0.0403} \\
\bottomrule
\end{tabular}
\end{table}

구간은 전부 \texttt{lab.pairboot.cluster\_boot} \textbf{BCa}(B$=$2000, 군집$=$격자) $\cdot$ \textbf{폴백 0 건}.

\claimx{\textbf{판정 §6-나 --- 겹친다.} ③ 의 첫째 자를 $0.5 \to 0$ 으로 되돌린다.}
{병기 판 셋은 전부 안 겹친다. 그래도 \textbf{사전등록이 미리 지정한 [둘 다] 판 하나로만} 판정했다 ---
좋게 나온 판으로 갈아타지 않았다(노트 133).}

\subsection{🔴 그리고 관문이 스스로 검정력 0 을 신고했다}

이번엔 모든 관문에 \textbf{도달 가능 폭 ÷ 문턱}을 붙였다.
주 판정 관문 G1 의 비가 \textbf{0.532}다 --- \textbf{1 미만이므로 검정력 0} 이 자동으로 찍혔다.

\claimx{그러므로 「겹친다」는 \textbf{「빈틈이 0.04025 일 미만이다」의 다른 말}이다.}
{이 표본은 그보다 작은 빈틈을 원리상 못 가른다. \textbf{「0 이다」라고 쓰면 그것은 과한 문장이다.}}
직전 노트의 관문 G2 도 다시 재니 비가 \textbf{0.0946} 이었다 --- 그것도 \textbf{검정력 0} 이고,
그래서 비교가능성의 정본 근거를 \textbf{홀드아웃 간격 다중집합 비트 동일}(54{,}862 행 전수)로 옮겼다.

\section{🔴 바닥에서 격자를 빼자 네 판이 모두 갈렸다}

바닥 팔이 \texttt{(dow, quarter)} 만 보는 「격자를 모르는 팔」이라,
\textbf{시간 구조가 전혀 없는 N2 세계에서도} 이 자는 원판에서 $+0.2827$ 로 유의하게 이긴다.
그 이득의 \textbf{92.8\%가 \texttt{prevmed}}(그 격자의 과거 간격 중앙값 $=$ 발생률 추정)다.

바닥에 \texttt{prevmed} 를 넣은 자$'$ 로 다시 재면:

\begin{table}[h]\centering\small
\begin{tabular}{lrr}
\toprule
절제 & 자(기후값 바닥) 빈틈 & 자$'$(격자를 아는 바닥) 빈틈 \\
\midrule
원판 & $+0.3122$ & $+0.1337$ \\
첫간격제외 & $+0.1436$ & $+0.1224$ \\
달력제거 & $+0.0404$ & $+0.0153$ \\
\textbf{둘 다} & $\mathbf{-0.0214}$ (겹침) & $\mathbf{+0.0020}$ (\textbf{안 겹침}) \\
\bottomrule
\end{tabular}
\end{table}

\claimx{그리고 [둘 다] 판에서 \textbf{귀무가 처음으로 0 을 문다} --- N2 $0.00350\,[-0.0016, +0.0082]$.}
{직전 노트의 사전등록이 「깨끗한 갈래」로 적어 두고 한 번도 발화하지 않았던 조건이다.
바닥을 고치니 발화했다.}

\claimx{\textbf{그러나 주 판정은 안 바꿨다.}}
{사전등록 §3 이 [둘 다] 판과 기후값 바닥을 미리 지정했다.
병기가 더 좋다고 갈아타면 그것이 정확히 노트 133 이 금지하는 것이다.}
그리고 크기를 정직하게 적는다 --- 그 개선은 \textbf{23.9 분}이다.

\section{가장자리 절과 구간 덮음}

\begin{figure}[h]\centering
\includegraphics[width=\linewidth]{figs/fig2.pdf}
\caption{(a) 홀드아웃의 4.150\%가 빈틈의 50.9\%를 만든다. (b) 🔴 900 노트 만에 처음으로 둘째 자로 판정했다.
(c) 🔴 검정력에 분모를 둘 적었다.}
\end{figure}

\subsection{가장자리 --- 그리고 못 가른 것}

\texttt{prevmed} 결측 \textbf{2{,}277 행}은 홀드아웃의 \textbf{4.150\%}이고
\textbf{홀드아웃 격자 수와 값이 같다}(격자마다 첫 간격이 정확히 하나) ---
🔴 \textbf{값이 같을 뿐 다른 분모다}(조항 60).
그 행의 평균 간격이 $58.540$ 일로 나머지 $13.732$ 일의 4.26 배이고, 빈틈의 \textbf{50.9\%}를 만든다.

\claimx{\textbf{가장자리다.}}
{🔴 다만 「가장자리라서 길다」와 「가장자리에서 시작한 격자가 원래 드문 격자다」를
\textbf{이 표로는 못 갈랐다} --- 못 갈랐다고 적는다.}

\subsection{🔴 900 노트 만에 처음으로 둘째 자로 판정했다}

③ 의 자는 \textbf{둘}이다 --- 「간격 예측이 기후값을 이기나」와 \textbf{「구간 덮음」}.
둘째 자로 판정한 사이클이 900 노트 중 \textbf{0} 이었다.

눈금을 \textbf{학습 안에서만} 맞췄다(눈금용 4{,}980 / 확인 1{,}215 격자 $\cdot$ 홀드아웃 겹침 \textbf{0}).
홀드아웃 덮음이 \textbf{조건부 0.7893 $\cdot$ 기후값 0.8089} 로,
눈금 전의 과대 덮음 $0.9047/0.8796$ 이 \textbf{둘 다 명목 $0.80$ 의 $\pm0.05$ 안}에 들어왔다.

\claimx{\textbf{그런데 둘째 조건을 못 넘었다} --- 조건부 폭 $26.194$ 일 $>$ 기후값 폭 $24.604$ 일.}
{덮음을 맞춘 자리에서도 $26.901$ 대 $24.267$ 이다.
\textbf{「안 했다」가 아니라 「했는데 반만 넘었다」}이다.}

\section{배선 --- 검정력에 분모를 둘 적었다}

심은 6 $\cdot$ 발화 6 $\cdot$ 놓친 0. 그런데 이번엔 \textbf{못 심은 3 도 분모에 넣었다}:
\textbf{6/6 $=$ 1.000}(심은 것 기준) 대 \textbf{6/9 $=$ 0.667}(관문 전체).

신설한 검사 중 하나가 직전 노트의 자를 무너뜨렸다:

\claimx{\textbf{항등 순열을 심으니 다중집합 검사가 그대로 통과시켰다} --- 눈이 멀다.}
{항등 판에서 N2 개선이 진짜와 \textbf{비트로 같아지고}($+0.7545204330866538$) 판정이 뒤집힌다.
그런데 직전 노트가 비교가능성의 근거로 삼은 다중집합 검사는 \textbf{그 판을 통과시킨다}.}

\section{한계 --- 「없다」가 아니라 「못 했다」}

\begin{itemize}
\item ⓪ 의 예상이 하나 틀렸다 --- 첫간격제외 판 군집은 2{,}277 이 아니라 \textbf{2{,}102} 다
      (홀드아웃 간격이 하나뿐인 격자 175 개가 통째로 빠진다). 결론에는 보수적인 방향
\item 🔴 \textbf{러너를 두 번 완주했다.} 1 차 산출물을 sha 와 함께 보관했고
      \textbf{공통 키 937 개 중 값이 다른 키 0 개}다. 더한 것은 병기 절$\cdot$진단 곡선$\cdot$이 자백뿐이고
      팔$\cdot$정의$\cdot$문턱$\cdot$판정 분기는 안 건드렸다
\item 「달력제거」가 달력의 몫만 떼는지 \textbf{못 갈랐다} --- 셀 해상도 저하와 뒤엉킨다
\item 직전 노트들의 모형$\cdot$부트 코드를 그대로 \texttt{import} 하므로 \textbf{그 관이 틀렸으면 원리상 못 잡는다}
\item 자$'$ 를 \textbf{주 판정으로 못 걸었다} --- 사전등록이 병기로 지정했다.
      \textbf{다음 사이클이 그 자로 사전등록해야 한다}
\item 조건부 구간이 왜 더 넓은지 \textbf{못 갈랐다}
\item 탐색 팔을 \textbf{안 돌렸다}
\end{itemize}

\end{document}
